%%%%%%%%%%  scan idx 112  |  folio 101  %%%%%%%%%%

\medskip
\noindent\textbf{Remark 4.6(a).}
Let $X_0$ be affine. A formal versal deformation of $X_0$ exists if
and only if the support of $\underline D^1_{X_0}$ is a finite set.
This does not imply, however, that all nonsmooth points of $X_0$ are
isolated: there are affine schemes $X_0$ over $k$ having a
nonisolated nonsmooth point for which
$\operatorname{Supp}\underline D^1_{X_0}$ is finite.

\medskip
\noindent\textbf{Remark 4.6(b).}
Henceforth, a deformation of $X_0$ means an object of
$\underline M_{X_0}$. Where confusion is possible, such an object
will be called an infinitesimal deformation of $X_0$.

\medskip
\noindent\textbf{Remark 4.6(c).}
Let $X_0$ be a scheme over $k$. We have fixed a complete noetherian
local ring $\Lambda$ with residue field $k$ and considered
deformations of $X_0$ over $\underline C_\Lambda$.

%%%%%%%%%%  scan idx 113  |  folio 102  %%%%%%%%%%

In practice one usually takes $\Lambda=k$, or $\Lambda=W(k)$ when
$k$ is a perfect field of positive characteristic. The inclusion
$\underline C_k\subset\underline C_\Lambda$ is fully faithful, and
$\underline M_{X_0}|_{\underline C_k}$ is simply the category of
deformations of $X_0$ over the objects of $\underline C_k$.
Consequently, if $(R,\mathfrak X)$ is a formal versal deformation of
$X_0$ over $\underline C_\Lambda$, then
\[
 (k\otimes_\Lambda R,\;k\otimes_\Lambda\mathfrak X)
\]
is a formal versal deformation of $X_0$ over $\underline C_k$. It
follows, for example, that
\[
 \dim_{\underline C_\Lambda}\underline M_{X_0}
 =\dim(k\otimes_\Lambda R)
\]
does not depend on the choice of $\Lambda$, but only on $X_0$.

The preceding theorem admits many generalizations. For example,
consider a pair $Y_0\longrightarrow X_0$, where $X_0$ is a scheme of
finite type over $k$ and $Y_0$ is a fixed closed subscheme of $X_0$.
For $R\in\operatorname{Ob}(\underline C_\Lambda)$, a deformation of
$(X_0,Y_0)$ over $R$ is a commutative diagram
\[
\begin{tikzcd}[row sep=large,column sep=huge]
Y_0 \arrow[rrr] \arrow[dr,hook] \arrow[ddr]
&&& Y \arrow[dl,hook] \arrow[ddl,"\mathrm{flat}"] \\
& X_0 \arrow[r] \arrow[d]
& X \arrow[d,"\mathrm{flat}"'] & \\
& \operatorname{Spec}(k) \arrow[r]
& \operatorname{Spec}(R) &
\end{tikzcd}
\]
such that the induced map
\[
(X_0,Y_0)\longrightarrow
\left(
 X\times_{\operatorname{Spec}(R)}\operatorname{Spec}(k),
 Y\times_{\operatorname{Spec}(R)}\operatorname{Spec}(k)
\right)
\]
is an isomorphism. Defining morphisms
\[
(R';X',Y')\longrightarrow(R;X,Y)
\]
of deformations of $(X_0,Y_0)$ in the evident way gives a cofibered
category $\underline M_{X_0,Y_0}$ in groupoids over
$\underline C_\Lambda$. This construction is frequently used to
rigidify automorphisms in deformation-moduli problems. The argument
of Corollary 4.3, based on Lemma 4.2, immediately shows that
$\underline M_{X_0,Y_0}$ is a homogeneous cofibered category in
groupoids over $\underline C_\Lambda$.
